\section{From prelocalized schobers to localized schobers}\label{section:pretolocalized}
In this section, we describe how we may recover localized schobers from prelocalized ones by stabilizing an operation on $\olsftwoP^{\pre}(\bbC,R)$ called symplectic suspension. Together with our study of prelocalized schobers in \S\ref{section:prelocalized}, this gives a precise sense in which we may define localized schobers from perverse schobers, as depicted below.
\begin{align*}
    \text{Schobers on }(\bbC,R)&\overset{\text{quotient}}{\rightsquigarrow}\text{Prelocalized schobers on }(\bbC,R)\\
    &\overset{\text{stabilization}}{\rightsquigarrow}\text{Localized schobers on }(\bbC,R)
\end{align*}
The need to stabilize the symplectic suspension is motivated by the following issue which arises in symplectic geometry.

\begin{obs}\label{obs:stable}
    Let $f:X\to\bbC$ be a Lefschetz (or symplectic) fibration with critical values $R$. Let $\calF(X,f)\in\sfP(\bbC,R)$ be the associated perverse schober defined in \cite{kapranovPerverseSchobersAlgebra2020}*{\S 3.6} and let $\ol{\calF(X,f)}^{\opname{pre}}\in\olsftwoP^{\opname{pre}}(\bbC,R)$ be the resulting prelocalized perverse schober.

    The operation $\overline{\calF(X,f)}^{\opname{pre}}\mapsto \overline{\calF(X\times \bbC_x,f+x^2)}^{\opname{pre}}$ of prelocalized schobers is not ``invertible''.
\end{obs}

\begin{eg}
    Consider for example $X=\bbC\xto{z^2}\bbC$. Then, $\calF(\bbC,z^2)$ is equivalent to the spherical adjunction
    \[\Mod_k\tofrom\Loc(S^0)\simeq\Mod_k\oplus\Mod_k.\]
    Thus the prelocalized schober is given by $\Mod_k$ together with the monad/algebra $C^\bullet(S^0;k)$ of cochains on $S^0$. The algebra structure is that of the cup product.

    After symplectic suspension, $\calF(\bbC^2,x^2+z^2)$ is equivalent to the spherical adjunction
    \[\Mod_k\tofrom\Loc(S^1).\]
    Thus the prelocalized schober is $\Mod_k$ with monad/algebra $C^\bullet(S^1;k)$.

    The algebra $C^\bullet(S^1;k)$ contains ``less information'' since this is isomorphic to a square-zero extension of $k$ while $C^\bullet(S^0;k)$ is not.
\end{eg}

In \S\ref{subsection:suspension}, we formulate the symplectic suspension operation algebraically, working at the level of individual abstract perverse schobers. We expect that under this operation, $\calF(X,f)\mapsto \calF(X\times \bbC_x,f+x^2)$, although we don't verify this outside examples such as the one above. We then define the symplectic suspension similarly for abstract prelocalized and localized perverse schobers.

We then upgrade these to functors of $\two$-categories and exhibit compatibility with braid group actions so that we obtain symplectic suspension functors on (prelocalized) perverse schobers,
\begin{align*}
    \calSus_{\sftwoP}:&\sftwoP(\bbC,R)\to \sftwoP(\bbC,R),\\
    \calSus_{\olsftwoP}:&\olsftwoP^{\pre}(\bbC,R)\to \olsftwoP^{\pre}(\bbC,R).
\end{align*}

We prove in \S\ref{subsection:stabilization} the main theorem of this section, which confirms that localized schobers can indeed be recovered from prelocalized schobers by stabilizing symplectic suspension.

\begin{thm}\label{thm:limit0}
    There is an equivalence of $\two$-categories,
    \[\olsftwoP(\bbC,R)\isomto \lim(\cdots\to\olsftwoP^{\pre}(\bbC,R)\xto{\calSus}\olsftwoP^{\pre}(\bbC,R)\xto{\calSus}\olsftwoP^{\pre}(\bbC,R)).\]
\end{thm}

\begin{proof}
    This follows from Theorem~\ref{thm:limit1}.
\end{proof}

\subsection{Symplectic suspension}\label{subsection:suspension}
We will begin by defining the symplectic suspension operation on $\sfAdj$ and show that it induces operations on $\sfSph$, $\sfMnd$ and $\sfSphMnd$. We will later upgrade these to functors of $\two$-categories.

On $\sfAdj$, consider the following operation defined on objects.
\begin{align*}
    \calSus_{\sfAdj}: \sfAdj&\to \sfAdj\\
    [S:\Phi\tofrom\Psi:R]&\mapsto [S':\Phi\tofrom\Psi':=\lim(\Phi\oplus\Phi\to\Psi):R'].
\end{align*}
Here, the limit is understood to be strict. So an object is given by a triple $(\phi_1,\phi_2,\alpha)$, where $\alpha:S\phi_1\isomto S\phi_2$ is an isomorphism. In this notation, the adjoint functors are given by
\begin{align*}
    S':\phi&\mapsto (\phi,\phi,1),\\
    R':(\phi_1,\phi_2,\alpha) &\mapsto \coCone(\phi_1\oplus\phi_2\to RS\phi_1).
\end{align*}
We calculate in Lemma~\ref{lem:sussph} that there is an identification $\coCone(\id\to R'S')\simeq \coCone(\id\to RS)[-1]$.

\begin{eg}\label{eg:localsystem}
    Let $X$ be a space and consider the adjunction
    \[\calA(X):=[\pi^*:\Mod_k\tofrom\Loc(X):\pi_*],\]
    where $\Loc(X)$ denotes the category of $\infty$-local systems. Then,
    \[\calSus_{\sfAdj}(\calA(X))\simeq\calA(\Sigma X).\]
    In this way, suspension of spaces is compatible with suspension of adjoint pairs. 
    %More generally, if $f:X\to Y$ is a map of spaces, and $\calA(f)$ is the adjoint pair $f^*\adjto f_*$, we have
    %\[\calSus(\calA(f))\simeq\calA(Y\sqcup_XY\to Y).\]
\end{eg}

\begin{lem}\label{lem:sussph}
    Suppose that an adjunction $\calA$ is spherical. Then its suspension $\calSus_{\sfAdj}(\calA)$ is also spherical. Thus the operation descends to $\calSus_{\sfSph}:\sfSph\to\sfSph$.
\end{lem}

\begin{proof}
    We may use the formulae above to compute the twist functors explicitly. We have
    \[T'_{\Phi}:\phi\mapsto \Tot(\phi\xto{(1,-1)}\phi\oplus\phi\to RS\phi)\simeq T_{\Phi}[-1]\phi,\]
    and
    \[T'_{\Psi'}:(\phi_1,\phi_2,\alpha)\mapsto (T_{\Phi}[1]\phi_2,T_{\Phi}[1]\phi_1,T_{\Psi}\alpha^{-1}).\]
    These are clearly autoequivalences.
\end{proof}

On $\sfMnd$, consider the following operation defined on objects.
\begin{align*}
    \calSus_{\sfMnd}: \sfMnd&\to \sfMnd\\
    (\Phi,M)&\mapsto (\Phi,M':=\lim(\id\oplus \id\to M))
\end{align*}
Here, $\id$ is the trivial monad on $\Phi$, $\id\to M$ is the canonical map of monads and the limit is taken in the category of monads on $\Phi$.

\begin{lem}\label{lem:sussphmnd}
    Suppose that a monad $M$ on $\Phi$ is spherical. Then its suspension $\calSus_{\sfMnd}(M)$ is also spherical. Thus the operation descends to $\calSus_{\sfSphMnd}:\sfSphMnd\to \sfSphMnd$.
\end{lem}

\begin{proof}
    This follows by observing that $\coCone(\id\to M')\simeq \coCone(\id\to M)[-1]$, and applying Theorem~\ref{thm:christstrengthening1}.
\end{proof}

Finally, on $\sfEnd$ and $\sfAut$, we have the operations $\calSus_{\sfEnd},\calSus_{\sfAut}$, each given by the formula $(\Phi,T)\mapsto (\Phi,T[-1])$.

For general $n$, we define
\[\calSus_{\sfAut}:\sfAut(n)\to \sfAut(n)\]
by $(\frakS;\Phi,T)\mapsto (\frakS;\Phi,T[-1])$. We can also define
\begin{align*}
    \calSus_{\sfSph}:&\sfSph(n)\to \sfSph(n),\\
    \calSus_{\sfSphMnd}:&\sfSphMnd(n)\to \sfSphMnd(n)
\end{align*}
by applying $\sfAut(n)\times_{\sfAut}-$ to $\calSus_{\sfSph},\calSus_{\sfSphMnd}$ respectively.

\begin{rmk}[Thom--Sebastiani for schobers]
    We may generalize this construction to a binary operation,
    \begin{align*}
        -\ast-: \sfSph\times \sfSph&\to \sfSph\\
        [S_1:\Phi_1\tofrom\Psi_1:R_1],[S_2:\Phi_2\tofrom\Psi_2:R_2]&\mapsto [S_{12}:\Phi_{12}\tofrom\Psi_{12}:S_{12}^R],
    \end{align*}
    so that $\calSus_{\sfSph} = \calF(\bbC_z,z^2)\ast-$. In future work, we study this binary operation. We expect that on Lefschetz schobers (with critical value $R=\{0\}$), there is an identification,
    \[\calF(X,f)\ast\calF(Y,g)\simeq \calF(X\times Y,f\boxplus g),\]
    which can be thought of as a categorification of the Thom--Sebastiani formula. Moreover, we expect that under the 3d mirror symmetry of Gammage--Hilburn--Mazel-Gee \cite{gammagePerverseSchobers3d2023}, this operation corresponds to convolution monoidal product $\opname{2IndCoh}_{\mathbb{L}}(\bbA^1/\bbG_m)$ coming from the product monoidal structure on $\bbA^1/\bbG_m$.
\end{rmk}

We now upgrade the symplectic suspension to endofunctors of $\two$-categories. To do this, it is natural to work in certain enlarged $\two$-categories, where we allow for more 1-morphisms given by oplax natural transformations of diagrams. This is already apparent from our definitions at the level of objects. For example, to define the symplectic suspension on monads, we wrote down a diagram $\id\to M\from \id$ of monads on $\Phi$. The arrow $(\Phi,\id)\to (\Phi,M)$ doesn't appear in $\sfMnd$, as this is given by an oplax natural transformation of presheaves on $\frakmnd$.

We recall from \cite{haugsengLaxTransformationsAdjunctions2021} that $\sfAdj$ is equivalent to the $\two$-category $\sfFun(\frakadj,\sfSt_k)$ of diagrams on the walking adjunction $\frakadj$. Consider the functors of $\two$-categories,
\[\frakadj\from \frakmnd\from \frakend\from \mathfrak{pt},\]
where $\mathfrak{pt}$ denotes the trivial $\two$-category. We denote the maps from the trivial $\two$-category by $\mathfrak{pt}\xto{a} \frakadj,\mathfrak{pt}\xto{m}\frakmnd,\mathfrak{pt}\xto{e}\frakend$. It follows by Theorem~\ref{thm:laxfibrations} that there are left adjoints to $a^*,m^*,e^*$, denoted $a_!,m_!,e_!$. These send $\Phi\in\sfSt_k$ to the objects
\[[\Phi\tofrom \Phi]\in\sfAdj,\; (\Phi,\id)\in\sfMnd,\; (\Phi,\id)\in\sfEnd\]
respectively. It follows that the units of these adjunctions, $\Id\to a^*a_!, m^*m_!, e^*e_!$ are each equivalences, and moreover that the resulting lift of $a_!\Phi$ to $\sfAdj_{\oplax}\resto{\Phi}$ is an initial object. Similar statements hold for $\sfMnd_{\oplax}$ and $\sfEnd_{\oplax}$.

We define functors
\begin{align*}
    \calSus_{\sfAdj_{\oplax}}&:=a_!a^*\times_{\id}a_!a^*:\sfAdj_{\oplax}\to\sfAdj_{\oplax}\\
    \calSus_{\sfMnd_{\oplax}}&:=m_!m^*\times_{\id}m_!m^*:\sfMnd_{\oplax}\to\sfMnd_{\oplax}\\
    \calSus_{\sfEnd_{\oplax}}&:=e_!e^*\times_{\id}e_!e^*:\sfEnd_{\oplax}\to\sfEnd_{\oplax}.
\end{align*}

These are compatible with the pullback functors
\[\sfAdj_{\oplax}\to \sfMnd_{\oplax}\to \sfEnd_{\oplax},\]
because pullback, being a right adjoint, commutes with limits.

These functors restrict to the locally full subcategories $\sfAdj\subset\sfAdj_{\oplax}$, $\sfMnd$ and $\sfEnd$ respectively. By construction, these agree on the level of objects with the mappings described in the previous subsection. It follows that symplectic suspension restricts to endofunctors of $\sfSph,\sfSphMnd,\sfAut$ respectively.

For general $n$, define $\calSus_{\sfAut}$ on $\sfAut(n)$ by recalling that it is a full sub-$\two$-category,
\[\sfAut(n)\subset\sfSOD(n)\times_{\sfSt_k}\sfAut,\]
and seeing that $\Id\times\calSus_{\sfAut}$ preserves this sub-$\two$-category. We can define the symplectic suspension functor on $\sfSph(n)$ by recalling from Proposition~\ref{prop:activecartesian2} that
\[\sfSph(n)\simeq \sfAut(n)\times_{\sfAut}\sfSph,\]
and applying symplectic suspension to each of these three terms separately. We can do similarly for $\sfSphMnd(n)$. It easily follows that symplectic suspension naturally commutes with the braid group action $\Br_n$, and so we obtain symplectic suspension operations
\[\calSus_{\sftwoP},\calSus_{\olsftwoP^{\pre}},\calSus_{\olsftwoP}\]
on the respective (non-abstract) perverse schober categories.

\subsection{Stabilization}\label{subsection:stabilization}
We may now translate Observation~\ref{obs:stable} into algebraic terms on the level of objects.

\begin{obs}\label{obs:key}
    When $n\geq1$, the functor $\calSus_{\sfSphMnd}:\sfSphMnd(n)\to \sfSphMnd(n)$ is not an equivalence.
\end{obs}

For geometric reasons explained in \S\ref{subsection:radonissuspension}, this is problematic. The following example demonstrates that it is not an equivalence on $\sfSphMnd$.

\begin{eg}\label{eg:suspendtwice}
    Assume that $k$ is a discrete commutative ring. Monads on $\Mod_k$ are given by $k$-algebras. We claim that if $A$ is a $k$-algebra and $B=k\times_Ak$ is its suspension, then the unit map $k\to B$, viewed as a morphism in $\Mod_k$, admits a splitting $B\to k$. Indeed, we can simply choose one of the two projections to $k$.

    So to see that $\calSus$ is not an equivalence on $\sfSphMnd$, it suffices to exhibit any $k$-algebra $B$ which doesn't admit such a splitting (and gives a spherical monad). We can simply take $B=0$.
\end{eg}

We fix this by inverting $\calSus$ on $\sfSphMnd$ and more generally on $\sfSphMnd(n)$ in a universal manner. The following theorem says that this procedure recovers the $\two$-category of abstract \textit{localized} perverse schobers.

\begin{thm}\label{thm:limit1}
    There is an equivalence of $\two$-categories,
    \[\sfAut(n)\isomto \lim(\cdots\to\sfSphMnd(n)\xto{\calSus}\sfSphMnd(n)\xto{\calSus}\sfSphMnd(n)).\]
\end{thm}

\begin{rmk}\label{rmk:splitting}
    We saw in Example~\ref{eg:suspendtwice} that any $k$-algebra $B$ which doesn't split as a $k$-module doesn't have a preimage along $\calSus$. We may interpret Theorem~\ref{thm:limit1} as saying that the $k$-algebras which admit all repeated preimages along $\calSus$ are those that split as a square-zero extension, $B=k\oplus T$, where $T$ is a $k$-module for which $T\otimes_k-$ is invertible.

    The construction appearing in \cite{segalAllAutoequivalencesAre2020} uses that from any autoequivalence $T$, we may produce a square-zero monad $\id\oplus T[-1]$. We expect that this is exactly what $\sfAut\to \sfSphMnd$ does.
\end{rmk}

\comment{
    Motivated by this, we make the following definition.    

    \begin{defn}
        Let $(\Phi,F)\in \sfSphMnd$ be an object. A splitting of $(\Phi,F)$ is an autoequivalence $T$ of $\Phi$ together with an identification of monads
        \[F\simeq \id\oplus T[1],\]
        where $\id\oplus T[1]$ is given the structure of a square-zero monad.

        Likewise, a splitting of $(\frakS;\Phi,F)\in\sfSphMnd(n)$ is an autoequivalence of $T$ of $\Phi$ together with an identification of monads
        \[F\simeq \id\oplus T[1],\]
        where $\id\oplus T[1]$ is given the structure of a square-zero monad.
    \end{defn}

    Observe that if $(\frakS;\Phi,\id\oplus T[1])\in\sfSphMnd(n)$ is an object with a splitting, then $T$ is necessarily conjugate-compatible with $\frakS$, and every split object arises in this way. As such, we may think of $\sfAut(n)$ as the $\two$-category of split monads. The following makes this relation precise.

    \begin{thm}
        For any $(\frakS;\Phi,F)\in\sfSphMnd(n)$, the image under the operation $\calSus_{\sfSphMnd}^2(\calF)$ admits a splitting
        \[\calSus_{\sfSphMnd}^2(\frakS;\Phi,F)\simeq (\frakS;\Phi,\id\oplus T[-1]).\]
        In particular, the image of $\calSus_{\sfSphMnd}^2$ consists of the subset of equivalence classes of objects that admit a splitting, and moreover $\calSus_{\sfSphMnd}^2$ is an a bijection on this subset.
    \end{thm}

    \begin{proof}
        We prove an analogous statement for $\calSus_{\sfMnd}^2:\sfMnd\to \sfMnd$.

        We begin by stating a formula for $\calSus_{\sfAdj}^2$.
        \begin{align*}
            \calSus^2: \sfAdj&\to \sfAdj\\
            [S:\Phi\tofrom\Psi:R]&\mapsto [S'':\Phi\tofrom\Psi'':=\lim(\Phi\otimes (\Sigma S^0)\oplus\Psi\to\Psi\otimes (\Sigma S^0)):R''].
        \end{align*}
        Geometrically, $\Psi''$ parametrizes a $\Sigma S^0\simeq S^1$-family of objects in $\Phi$ together with a trivialization of this after the application of $S$. We choose say, the North pole of $\Sigma S^0\simeq S^1$ so that we may write
        \[\Psi''\simeq \{(\phi,u,\beta:\id_{S\phi}\simeq Su)\}.\]

        Let $F := 1\oplus T[-1]$ denote the desired square-free monad. Our strategy is to produce a functor $A:\Psi''\to\Mod_{\Phi}(F)$ together with a factorization as follows.
        \begin{equation}\label{equation:factor}
            \begin{tikzcd}
                {\Phi} & \Psi'' \\
                {} & \Mod_{\Phi}(F)
                \arrow["A", dashed, from=1-2, to=2-2]
                \arrow["R''"', from=1-2, to=1-1]
                \arrow["\opname{Forget}", from=2-2, to=1-1]
            \end{tikzcd}
        \end{equation}

        We construct $A$ explicitly as follows. Given object $\psi'' := (\phi,u,\beta)\in \Psi''$, the underlying object of our $F$-module is $\coCone(\phi\xto{v} T\phi)$. Here, the map denoted $v$ comes from the following diagram in $\Phi$ where the bottom row is an exact triangle, and the commuting triangle comes from $\beta$.
        \[\begin{tikzcd}
            & \phi & & \\
            T\phi & \phi & RS\phi & {}
            \arrow["v"{description}, dashed, from=1-2, to=2-1]
            \arrow["1-u"{description}, from=1-2, to=2-2]
            \arrow["0"{description}, from=1-2, to=2-3]
            \arrow[from=2-1, to=2-2]
            \arrow[from=2-2, to=2-3]
            \arrow["+1", from=2-3, to=2-4]
        \end{tikzcd}\]
        With some care, this assignment upgrades to a functor $\Psi''\to \Phi$.

        To write down the action of $F$ on $A\psi''$, we first write down a map $a:T[-1](A\psi'')\to A\psi''$. We define this via the composition,
        \[T[-1]\coCone(\phi\xto{v} T\phi)\to T[-1]\phi \to \coCone(\phi\xto{v} T\phi).\]
        Now the associativity constraints and higher constraints can be easily constructed. For example, we must exhibit an identification of
        \[(T[-1])^2(A\psi'')\xto{T[-1]a} (T[-1])(A\psi'')\xto{a} A\psi''\]
        with zero. There is a canonical one by observing that when we expand this into the composition of four maps, the middle two are
        \[(T[-1])^2\phi \to T[-1]\coCone(\phi\xto{v} T\phi)\to T[-1]\phi.\]
        Higher associativity constraints are similarly canonically constructed.

        We next show that the map of functors $R''S''\to F$ is an equivalence. We take left adjoints in (\ref{equation:factor}) to get the following lax-commuting diagram.
        \begin{equation}\label{equation:factor2}
            % https://q.uiver.app/#q=WzAsMyxbMCwwLCJcXFBzaScnIl0sWzEsMCwiXFxQaGkiXSxbMCwxLCJcXE1vZF97XFxQaGl9KEYpIl0sWzEsMCwiUycnIiwyXSxbMCwyLCJBIiwyXSxbMSwyLCJcXG9wbmFtZXtGcmVlfSJdLFs1LDAsIiIsMCx7InNob3J0ZW4iOnsic291cmNlIjoyMH19XV0=
            \begin{tikzcd}[ampersand replacement=\&]
                {\Psi''} \& \Phi \\
                {\Mod_{\Phi}(F)}
                \arrow["A"', from=1-1, to=2-1]
                \arrow["{S''}"', from=1-2, to=1-1]
                \arrow[""{name=0, anchor=center, inner sep=0}, "{\opname{Free}}", from=1-2, to=2-1]
                \arrow[between={0.2}{1}, Rightarrow, from=0, to=1-1]
            \end{tikzcd}
        \end{equation}
        Now, $S''\phi = (\phi,\id_{\phi},\id_{\id_{\phi}})$ produces under $A$, the object
        \[\coCone(\phi\xto{0}T\phi) = \phi \oplus T[-1]\phi\]
        and the lax-commutativity
        \[\opname{Free}(\phi) \simeq \phi\oplus T[-1]\phi \to A(S''\phi).\]
        is in fact strict. This shows that indeed, $R''S''\isomto F$ is an equivalence. It follows that the monad induced by $S''\adjto R''$ is identified with $F$, completing the proof.
    \end{proof}
}


\begin{prop}\label{prop:limit2}
    Consider the following diagram in which each row is a limit diagram, and where we have omitted the subscript of $\calSus$ for readability.
    % https://q.uiver.app/#q=WzAsMTAsWzIsMCwiXFxzZk1uZCJdLFszLDAsIlxcc2ZNbmQiXSxbNCwwLCJcXHNmTW5kIl0sWzIsMSwiXFxzZkVuZCJdLFszLDEsIlxcc2ZFbmQiXSxbNCwxLCJcXHNmRW5kIl0sWzEsMCwiXFxjZG90cyJdLFsxLDEsIlxcY2RvdHMiXSxbMCwwLCJcXHNmTW5kW1xcY2FsU3VzXnstMX1dIl0sWzAsMSwiXFxzZkVuZCJdLFswLDEsIlxcY2FsU3VzIl0sWzEsMiwiXFxjYWxTdXMiXSxbMyw0LCJcXGNhbFN1cyIsMl0sWzQsNSwiXFxjYWxTdXMiLDJdLFswLDNdLFsxLDRdLFsyLDVdLFs2LDBdLFs3LDNdLFs5LDddLFs4LDZdLFs4LDldXQ==
    \[\begin{tikzcd}[ampersand replacement=\&]
        {\sfMnd[\calSus^{-1}]} \& \cdots \& \sfMnd \& \sfMnd \& \sfMnd \\
        \sfEnd \& \cdots \& \sfEnd \& \sfEnd \& \sfEnd
        \arrow[from=1-1, to=1-2]
        \arrow["\opname{C}"', from=1-1, to=2-1]
        \arrow[from=1-2, to=1-3]
        \arrow["\calSus", from=1-3, to=1-4]
        \arrow[from=1-3, to=2-3]
        \arrow["\calSus", from=1-4, to=1-5]
        \arrow[from=1-4, to=2-4]
        \arrow[from=1-5, to=2-5]
        \arrow[from=2-1, to=2-2]
        \arrow[from=2-2, to=2-3]
        \arrow["\calSus"', from=2-3, to=2-4]
        \arrow["\calSus"', from=2-4, to=2-5]
    \end{tikzcd}\]
    The resulting functor $\opname{C}:\sfMnd[\calSus^{-1}]\to \sfEnd$ is an equivalence of $\two$-categories. 
\end{prop}

\begin{proof}
    By Yoneda lemma, it suffices to show that for any test object $\frakt\in\wCat_{\two}$, the following map of mapping spaces is an equivalence,
    \begin{equation}\label{equation:wantequiv}
        \Map_{\wCat_{\two}}(\frakt,\sfMnd[\calSus^{-1}])\to \Map_{\wCat_{\two}}(\frakt,\sfEnd).
    \end{equation}
    In fact, this lives over the $\zero$-category $\Map(\frakt,\sfSt_k)$. It suffices to check that for each point $\opname{X}\in \Map(\frakt,\sfSt_k)$, the map (\ref{equation:wantequiv}) restricts to an equivalence.

    We apply $\Map_{\wCat_{\two}}(\frakt,-)$ and then the restriction to the fiber over $\opname{X}$ to the whole diagram, noting that these operations preserve limits. By Lemma~\ref{lem:haugsengvariant}, we have a functorial identification,
    \[\Map(\frakt,\sfMnd)\resto{\opname{X}}\simeq \Map(\frakmnd,\sfFun(\frakt,\sfSt_k))\resto{\opname{X}}\simeq \Alg(\sfT(\opname{X},\opname{X}))^{\simeq},\]
    where $\sfT := \sfFun(\frakt,\sfSt_k)$. Similarly,
    \[\Map(\frakt,\sfEnd)\resto{\opname{X}}\simeq \Map(\frakmnd,\sfFun(\frakt,\sfSt_k))\resto{\opname{X}}\simeq \sfT(\opname{X},\opname{X})^{\simeq}.\]
    This results in the following diagram.
    % https://q.uiver.app/#q=WzAsOCxbMiwwLCJcXEFsZyhcXEVuZF97XFxzZlh9KFxcb3BuYW1le1h9KSlee1xcc2ltZXF9Il0sWzMsMCwiXFxBbGcoXFxFbmRfe1xcc2ZYfShcXG9wbmFtZXtYfSkpXntcXHNpbWVxfSJdLFsyLDEsIlxcRW5kX3tcXHNmWH0oXFxvcG5hbWV7WH0pXntcXHNpbWVxfSJdLFszLDEsIlxcRW5kX3tcXHNmWH0oXFxvcG5hbWV7WH0pXntcXHNpbWVxfSJdLFswLDAsIlxcbGltKFxcQWxnKFxcRW5kX3tcXHNmWH0oXFxvcG5hbWV7WH0pKV57XFxzaW1lcX0pIl0sWzEsMCwiXFxjZG90cyJdLFswLDEsIlxcRW5kX3tcXHNmWH0oXFxvcG5hbWV7WH0pXntcXHNpbWVxfSJdLFsxLDEsIlxcY2RvdHMiXSxbMCwxLCJcXGNhbFN1cyJdLFsyLDMsIlstMV0iLDJdLFswLDJdLFsxLDNdLFs0LDVdLFs1LDBdLFs2LDddLFs0LDZdLFs3LDJdXQ==
    \[\begin{tikzcd}[ampersand replacement=\&]
        {\lim(\Alg(\sfT(\opname{X},\opname{X}))^{\simeq})} \& \cdots \& {\Alg(\sfT(\opname{X},\opname{X}))^{\simeq}} \& {\Alg(\sfT(\opname{X},\opname{X}))^{\simeq}} \\
        {\sfT(\opname{X},\opname{X})^{\simeq}} \& \cdots \& {\sfT(\opname{X},\opname{X})^{\simeq}} \& {\sfT(\opname{X},\opname{X})^{\simeq}}
        \arrow[from=1-1, to=1-2]
        \arrow[from=1-1, to=2-1]
        \arrow[from=1-2, to=1-3]
        \arrow["\calSus", from=1-3, to=1-4]
        \arrow[from=1-3, to=2-3]
        \arrow[from=1-4, to=2-4]
        \arrow[from=2-1, to=2-2]
        \arrow[from=2-2, to=2-3]
        \arrow["{[-1]}"', from=2-3, to=2-4]
    \end{tikzcd}\]
    The left vertical arrow is an equivalence by applying the limit-preserving functor $(-)^{\simeq}$ to Lemma~\ref{lem:lurievariant} and identifying the map with the equivalence
    \[\Sp(\Alg(\sfT(\opname{X},\opname{X})))\simeq \sfT(\opname{X},\opname{X}),\]
    of \cite{lurieHigherAlgebra2017}*{7.3.4.7}.
\end{proof}

\begin{lem}\label{lem:haugsengvariant}
    Let $\sfT\in\wCat_{\two}$ and $\opname{X}\in\sfT$. Then, the fiber over $\opname{X}$ of the map of $\zero$-categories, $\Map(\frakmnd,\sfT)\to \sfT^{\simeq}$, is the $\zero$-category $\Alg(\sfT(\opname{X},\opname{X}))^{\simeq}$.  An analogous statement holds with $\frakmnd$ replaced by $\frakend$.
\end{lem}

\begin{proof}
    It follows from \cite{haugsengLaxTransformationsAdjunctions2021}*{Theorem 1.4} that the fiber of
    \[\sfFun(\frakmnd,\sfT)_{\oplax}^{\one}\to \sfT^{\one}\]
    over $\opname{X}$ is identified with $\Alg(\sfT(\opname{X},\opname{X}))$. The inclusion
    \[\sfFun(\frakmnd,\sfT)^{\one}\to \sfFun(\frakmnd,\sfT)_{\oplax}^{\one}\]
    induces an equivalence of $\zero$-categories after applying $(-)^{\simeq}$. This is because any oplax natural transformation of diagrams that yields an invertible 1-morphism in $\sfFun(\frakmnd,\sfT)_{\oplax}^{\one}$ is in fact an honest natural transformation of diagrams. This completes the proof. The statement for $\frakend$ is similar.
\end{proof}

\begin{lem}\label{lem:lurievariant}
    Let $\calC$ be a large, cocomplete $\one$-category with an initial object denoted $k\in\calC$. Let $\Omega_{\calC}$ be the operation, $c\mapsto k\times_c k$. Then the following limit exists, and there is an equivalence of categories,
    \[\opname{Sp}(\calC_{/k})\isomto\lim(\cdots\to \calC\xto{\Omega_{\calC}}\calC\xto{\Omega_{\calC}}\calC).\]
    Here, $\opname{Sp}(\calC_{/k})$ is the $\one$-category of spectrum objects of the pointed category $\calC_{/k}$ as defined in \cite{lurieHigherAlgebra2017}.
\end{lem}

\begin{rmk}
    This notation is confusing as our symplectic \textit{suspension} corresponds to a ``\textit{loop space} functor'' here. We also see this in the negative shift $[-1]$ in $\calSus_{\sfAut}$. This comes down to the covariance of the functor
    \[C^*(-;k):\opname{Spaces}\to \Alg_k.\]
\end{rmk}

\begin{proof}[Proof of Lemma~\ref{lem:lurievariant}]
    We observe that $\Omega_{\calC}$ factors through the slice category $\calC_{/k}$ once we choose one of the two projections, $\Omega_{\calC}c\to k$.
    \[\begin{tikzcd}
        \calC & \calC \\
        & {\calC_{/k}}
        \arrow["{\Omega_{\calC}}", from=1-1, to=1-2]
        \arrow[from=1-1, to=2-2]
        \arrow["{\opname{forget}}"', from=2-2, to=1-2]
    \end{tikzcd}\]
    We may extend this to the following naturally commuting diagram, where $\Omega_{\calC_{/k}}$ denotes the operation given by the same formula on the pointed category $\calC_{/k}$.
    % https://q.uiver.app/#q=WzAsOSxbMywwLCJcXGNhbEMiXSxbNCwwLCJcXGNhbEMiXSxbNCwxLCJcXGNhbENfey9rfSJdLFsyLDAsIlxcY2FsQyJdLFszLDEsIlxcY2FsQ197L2t9Il0sWzEsMCwiXFxjYWxDIl0sWzIsMSwiXFxjYWxDX3sva30iXSxbMSwxLCJcXGNkb3RzIl0sWzAsMCwiXFxjZG90cyJdLFsyLDFdLFswLDJdLFswLDEsIlxcT21lZ2Ffe1xcY2FsQ30iXSxbMyw0XSxbNSw2XSxbMywwLCJcXE9tZWdhX3tcXGNhbEN9Il0sWzUsMywiXFxPbWVnYV97XFxjYWxDfSJdLFs0LDIsIlxcT21lZ2Ffe1xcY2FsQ197L2t9fSIsMl0sWzYsNCwiXFxPbWVnYV97XFxjYWxDX3sva319IiwyXSxbNCwwXSxbNiwzXSxbNyw2XSxbOCw1XV0=
    \[\begin{tikzcd}[ampersand replacement=\&]
        \cdots \& \calC \& \calC \& \calC \& \calC \\
        \& \cdots \& {\calC_{/k}} \& {\calC_{/k}} \& {\calC_{/k}}
        \arrow[from=1-1, to=1-2]
        \arrow["{\Omega_{\calC}}", from=1-2, to=1-3]
        \arrow[from=1-2, to=2-3]
        \arrow["{\Omega_{\calC}}", from=1-3, to=1-4]
        \arrow[from=1-3, to=2-4]
        \arrow["{\Omega_{\calC}}", from=1-4, to=1-5]
        \arrow[from=1-4, to=2-5]
        \arrow[from=2-2, to=2-3]
        \arrow[from=2-3, to=1-3]
        \arrow["{\Omega_{\calC_{/k}}}"', from=2-3, to=2-4]
        \arrow[from=2-4, to=1-4]
        \arrow["{\Omega_{\calC_{/k}}}"', from=2-4, to=2-5]
        \arrow[from=2-5, to=1-5]
    \end{tikzcd}\]
    The top and bottom rows are coinitial in each other. Thus the existence of the limit, and the value of the limit of these rows are identified. The limit of the bottom row exists, and is the definition of the stabilization $\opname{Sp}(\calC_{/k})$.
\end{proof}

\begin{proof}[Proof of Theorem~\ref{thm:limit1}]
    We first prove the statement when $n=1$. We can pullback the diagram of Proposition~\ref{prop:limit2} along $\sfAut\inclto \sfEnd$. Using Theorem~\ref{thm:christstrengthening1}, the top row of the resulting diagram, which is still a limit, is identified with
    \[\sfAut\to\cdots\to\sfSphMnd\xto{\calSus}\sfSphMnd\xto{\calSus}\sfSphMnd,\]
    exhibiting the desired limit. The result for general $n$ follows by pulling back along $\sfAut(n)\to\sfAut$.
\end{proof}